

\section{Задание 1. Непрерывное и дискретное преобразование Фурье}
Прежде, чем приступить к выполнению задания, рассмотрим знакомую прямоугольную функцию $\Pi : \mathbb{R} \to \mathbb{R}$:

$$
\Pi(t) = 
\begin{cases} 
1, & |t| \leq 1/2, \\
0, & |t| > 1/2
\end{cases}
$$

Её Фурье-образом будет являться аналитическое выражение:

$$
\hat{\Pi}(\nu) = \int_{-\infty}^{+\infty} \Pi(t) e^{-2\pi i\nu t}  dt
$$

\text{Вычислим вручную аналитически образ Фурье функции $\Pi(t)$:}\\
$$
\hat{\Pi}(\nu) = \int_{-\infty}^{+\infty} \Pi(t) e^{-2\pi i\nu t}  dt=\int_{-0.5}^{0.5}e^{-2\pi i\nu t}dt=\int_{-0.5}^{0.5}\frac{1}{-2\pi i\nu}de^{-2\pi i\nu t}=\frac{e^{-2\pi i\nu t}}{-2\pi i\nu} \Bigg|_{-0.5}^{0.5}=\frac{\sin(\pi\nu)}{\pi\nu}=\text{sinc}(\pi\nu)
$$
\text{Построим графики функций $\Pi(t)$ (см. рис. \ref{1}) и $\hat{\Pi}(\nu)$ (см. рис. \ref{2})}\\

\begin{figure}[H]
    \centering
    \includegraphics[width=0.72\linewidth]{/home/ns/Desktop/ITMO-courses/Fourier_Methods/Lab5/plots/П(t).png}
    \caption{График $\Pi(t)$}
    \label{1}
\end{figure}

\begin{figure}[H]
    \centering
    \includegraphics[width=0.72\linewidth]{/home/ns/Desktop/ITMO-courses/Fourier_Methods/Lab5/plots/П_hat(v).png}
    \caption{График $\hat{\Pi}(\nu)$}
    \label{2}
\end{figure}

\subsection{Численное интегрирование}
Найдём Фурье-образ функции $\Pi(t)$ с помощью численного интегрирования (функция \texttt{trapz}).\\
Метод трапеций для численного приближённого вычисления определённого интеграла от функции $f(x)$ на интервале $[a,b]$ выглядит так:
$$
S = \sum_{i=1}^{N-1}\frac{f(x_i)+f(x_{i-1})}{2}\cdot\Delta x,\quad\Delta x=x_i-x_{i-1},\quad N=\frac{b-a}{\Delta x}
$$
Применять метод будем для вычисления Фурье от наших функций ($-\infty$ и $\infty$ заменятся на пределы численого интегрирования):
$$
\begin{cases}
\Pi(t) = \int\limits_{-\infty}^{\infty} \hat{\Pi}(\nu) e^{2\pi i\nu t} d\nu \\
\hat{\Pi}(\nu) = \int\limits_{-\infty}^{\infty} \Pi(t) e^{-2\pi i\nu t} dt
\end{cases}
$$


$\ \ \ \ $Исследуем влияние параметров $T$ (промежуток интегрирования по времени) и $\Delta t$ (шаг дискретизации)
на получаемый путём численного интегрирования Фурье-образ функции (см. рис. \ref{3}) и восстановленную функцию из образа
(см. рис. \ref{4}), зафиксируем $V, \Delta\nu$.

Исследуем влияние параметров $V$ (промежуток интегрирования по времени) и $\Delta\nu$ (шаг дискретизации)
на получаемый путём численного интегрирования Фурье-образ функции (см. рис. \ref{5}) и восстановленную функцию из образа
(см. рис. \ref{6}), зафиксируем $T, \Delta t$.
\newline\newline
\textbf{Анализ и выводы}\\
Видим, что уменьшение шага интегрирования по времени даёт более точные результаты
при вычислении Фурье-образа, сильно увеличивается период в Фурье-образе, за счет чего пропадают 
наблюдаемые пики вблизи главных частот, и при восстановлении сигнал не имеет сильного искажения.
Увеличение периода $T$ приводит к появлению наблюдаемой периодичности в восстановленном сигнале, при этом не имеет большого
влияния на форму самого восстановленного сигнала.
Уменьшение шага интегрирования по частоте приводит к более точной аппроксимации Фурье-образа, но 
не оказывает влияния на форму восстановленного сигнала (если только не брать слишком большие значения, при которых вся важная информация просто потеряется).
При увеличении $V$ сигнал восстанавливается лучше, можем наблюдать эффект Гиббса.
При правильно подобранных параметрах $\Delta\nu,\Delta t, V, T$ метод работает с хорошей точностью, но вычисления занимают много времени.

\newpage
\begin{figure}[H]
    \centering
    \includegraphics[width=0.95\linewidth]{/home/ns/Desktop/ITMO-courses/Fourier_Methods/Lab5/plots/П_hat(v)_comparison_fix_v.png}
    \caption{Сравнение графиков $\hat{\Pi}(\nu)$}
    \label{3}
\end{figure}

\newpage
\begin{figure}[H]
    \centering
    \includegraphics[width=0.95\linewidth]{/home/ns/Desktop/ITMO-courses/Fourier_Methods/Lab5/plots/П(t)_comparison_fix_v.png}
    \caption{Сравнение графиков $\Pi(t)$}
    \label{4}
\end{figure}
\newpage
\begin{figure}[H]
    \centering
    \includegraphics[width=0.95\linewidth]{/home/ns/Desktop/ITMO-courses/Fourier_Methods/Lab5/plots/П_hat(v)_comparison_fix_t.png}
    \caption{Сравнение графиков $\hat{\Pi}(\nu)$}
    \label{5}
\end{figure}

\newpage
\begin{figure}[H]
    \centering
    \includegraphics[width=0.95\linewidth]{/home/ns/Desktop/ITMO-courses/Fourier_Methods/Lab5/plots/П(t)_comparison_fix_t.png}
    \caption{Сравнение графиков $\Pi(t)$}
    \label{6}
\end{figure}

\subsection{Использование DFT}
Найдём Фурье-образ функци $\Pi(t)$ с помощью дискретного преобразования Фурье, используя его так, чтобы преобразования было унитарным. Выполним обратное преобразование
от найденного Фурье-образа с помощью обратного дискретного преобразования.
Формулы прямого и обратного унитарного $DFT$:
$$
\begin{cases}
    X_m=\frac{1}{\sqrt{N}}\sum_{n=0}^{N-1}x_ne^{-i2\pi (\frac{mn}{N})} \\
    x_n=\frac{1}{\sqrt{N}}\sum_{n=0}^{N-1}X_me^{i2\pi \frac{mn}{N}}
\end{cases}
$$
Связь параметров $\Delta\nu$ c $T$, и $\Delta t$ с $V$ выглядит так:
$$
\Delta\nu=\frac{1}{T},\quad\Delta t=\frac{1}{V}
$$
Для вычислений будем использовать быстрое преобразование Фурье - алгоритм, позволяющий вычислять дискретное преобразование
Фурье за $O(N\log (N))$ действий.

Исследуем влияние параметров $T$ (промежуток интегрирования по времени) и $\Delta t$ (шаг дискретизации)
на получаемый при помощи fft Фурье-образ функции и восстановленную функцию из образа
(см. рис. \ref{7}-\ref{10}). Найдем соответствующие им параметры $V,\quad\Delta\nu$.

\textbf{Анализ и выводы}\\
Уменьшение шага дискретизации дает больший охват спектра частот, что даёт более точное описание сигнала, восстановления происходит точнее.
Увеличение периода $T$ также повышает точность восстановления сигнала благодаря тому, что уменьшается шаг дискретизации в частотной области и спектр
частот сигнала описывается более точно.\\
При этом, Фурье-образ не сходится с вычисленным аналитически, потому что дискретное преобразование Фурье предусмотрено для
дискретных непериодичных функций.\\
Так же при использовании DFT получается избежать эффекта Гиббса, и в целом восстановленный сигнал при правильно подобранных параметрах
намного лучше совпадает с исходным. Работает данный метод кратно быстрее численного интегрирования из предыдущего пункта.

\newpage
\begin{figure}[H]
    \centering
    \includegraphics[width=0.95\linewidth]{/home/ns/Desktop/ITMO-courses/Fourier_Methods/Lab5/plots/П_hat(v)_comparison_fft_fix_v.png}
    \caption{Сравнение графиков $\hat{\Pi}(\nu)$}
    \label{7}
\end{figure}

\newpage
\begin{figure}[H]
    \centering
    \includegraphics[width=0.95\linewidth]{/home/ns/Desktop/ITMO-courses/Fourier_Methods/Lab5/plots/П(t)_comparison_fft_fix_v.png}
    \caption{Сравнение графиков $\Pi(t)$}
    \label{8}
\end{figure}
\newpage
\begin{figure}[H]
    \centering
    \includegraphics[width=0.95\linewidth]{/home/ns/Desktop/ITMO-courses/Fourier_Methods/Lab5/plots/П_hat(v)_comparison_fft_fix_t.png}
    \caption{Сравнение графиков $\hat{\Pi}(\nu)$}
    \label{9}
\end{figure}

\newpage
\begin{figure}[H]
    \centering
    \includegraphics[width=0.95\linewidth]{/home/ns/Desktop/ITMO-courses/Fourier_Methods/Lab5/plots/П(t)_comparison_fft_fix_t.png}
    \caption{Сравнение графиков $\Pi(t)$}
    \label{10}
\end{figure}

\subsection{Объяснения}
\textbf{Проблема Фурье-образа}\\
Метод \texttt{trapz} фактически аппроксимирует вычисление аналитически найденного интеграла, поэтому при увеличении $T$ и устремлении $\Delta t$ к нулю
мы получаем истинный Фурье-образ (с некоторой погрешностью). DFT же должен работать с дискретными  и конечными, а не непрерывными функциями, поэтому при нахождении 
Фурье-образа от непрерывной фунции, заданной на бесконечности, мы не получим истинного Фурье-образа. Однако, в силу обратимости преобразований DFT и обратного
DFT, мы можем восстановить исходную функцию корректно.

\textbf{Точность обратного преобразования}\\
Метод \texttt{trapz} справляется с восстановлением исходной функции хуже DFT, потому что у прямоугольной функции есть скачки, из-за которых возникает эффект
Гиббса в восстановленной функции. DFT справляется с функцией благодаря обратимости преобразования.

\textbf{Скорость}\\
Для вычисления интеграла методом трапеций по $N$ точкам необходимо произвести суммирование $N$ раз. Соответственно, чтобы получить $N$ значений
функции, нам понадобится $N^2$ операций, время работы алгоритма в Big O нотации равно $O(N^2)$. FFT же является алгоритом для ускоренного вычисления DFT,
основанным на использовании симметрии матрицы преобразования Фурье, который работает за время $O(N\log(N))$. Наглядное сравнение времени работы алгоритмов
при увеличении количества данных представлено на рисунке \ref{11}.

\begin{figure}[H]
    \centering
    \includegraphics[width=0.95\linewidth]{/home/ns/Desktop/ITMO-courses/Fourier_Methods/Lab5/plots/N2_NlogN.png}
    \caption{Сравнение времени работы алгоритмов}
    \label{11}
\end{figure}

\subsection{Приближение непрерывного с помощью DFT}
Найдём способ получить правильный Фурье-образ, соответствующий непрерывному преобразованию Фурье, используя функцию \texttt{fft} и не прибегая к численному
интегрированию. Также найдем способ восстановить исходный сигнал по полученному Фурье-образу с помощью \texttt{ifft}.\\
Воспользуемся суммой Римана:
$$
\int_a^b f(t) dt \approx \sum_{n=0}^{N-1} f(t_n) \Delta t, \quad \text{где } t_n = a + n \Delta t, \quad \Delta t = \frac{b-a}{N-1}
$$
Аппроксимируем преобразование Фурье суммой Римана. Это можно сделать так:
$$
\hat{f}(\nu)=\int_{-T/2}^{T/2}f(t)e^{-2\pi i\nu t}dt\approx\sum_{n=0}^{N-1}f(t_n)e^{-2\pi i\nu t_n}\Delta t,\quad\text{где } t_n=-\frac{T}{2}+n\Delta t,
$$
$$
\Delta t=\frac{T}{N-1},\quad\nu=\frac{m}{N\Delta t}
$$
Или так:
$$
\hat{f}(\nu) \approx \int_{-T/2}^{T/2} f(t) e^{-2\pi i \nu t}  
dt \quad \rightarrow \quad \hat{f}(\nu_m) \approx c_m \sum_{n=0}^{N-1} f(t_n) e^{-2\pi i \frac{nm}{N}}
$$
Отсюда найдём коэффициент $c_m$:
$$
\sum_{n=0}^{N-1}f(t_n)e^{-2\pi i\nu t_n}\Delta t=c_m\sum_{n=0}^{N-1} f(t_n) e^{-2\pi i \frac{nm}{N}}
$$
$$
c_m=\Delta t \cdot e^{-2\pi i\left(\nu t_n-\frac{mn}{N}\right)}=
\Delta t \cdot e^{-2\pi i\left(\frac{m}{N\Delta t} \left(n\Delta t-\frac{T}{2}\right)-\frac{mn}{N}\right)}=
\Delta t \cdot e^{2\pi i\nu\frac{T}{2}}
$$
Или, если записать в более общем виде:
$$
c_m=\Delta t\cdot e^{-2\pi i\nu t_0},\quad\text{где }\nu=\frac{m}{N\Delta t}
$$
Теперь, чтобы получить Фурье-образ,найдём коэффициенты $c_m$ для всех частот и поэлементно домножим на массив точек, полученный из \texttt{fft} выполненного
для функции $\Pi(t)$. И найдем приближение исходной функции, поделив массив точек из \texttt{ifft} на тот же коэффициент. Посмотрим на графики и убедимся, что данный алгоритм сработал 
(\ref{12}).

\begin{figure}[H]
    \centering
    \includegraphics[width=0.95\linewidth]{/home/ns/Desktop/ITMO-courses/Fourier_Methods/Lab5/plots/П_hat(v)_comparison_sfft.png}
    \caption{Применение умного fft}
    \label{12}
\end{figure}
\begin{figure}[H]
    \centering
    \includegraphics[width=0.95\linewidth]{/home/ns/Desktop/ITMO-courses/Fourier_Methods/Lab5/plots/П(t)_comparison_sfft.png}
    \caption{Применение умного ifft}
    \label{13}
\end{figure}